\chapter{Contrats d'options et bourses d'options}\label{ch:m1:option-contracts-and-exchanges}

Une grande entreprise américaine n'a qu'une seule action. Sur cette action, ce matin, il existe
environ deux mille contrats d'options différents~: calls et puts, à soixante-dix
\omterm{def:m1:option-contracts-and-exchanges:option}{prix d'exercice} différents, pour des \omterm{def:m1:option-contracts-and-exchanges:option}{échéances} allant de ce vendredi à deux ans. Chacun des
deux mille est coté à double sens sur plus de quinze bourses à la fois,
par des \omterm{def:m1:what-a-trading-firm-does:market-maker}{teneurs de marché} qui doivent mettre à jour chacune de ces cotations à chaque
fois que l'action bouge d'un cent. Le marché des actions du \cref{ch:m1:us-equity-market-structure}
compte quelques milliers d'instruments~; celui des options, bâti sur lui, en
compte des centaines de fois davantage. Ce chapitre décrit le contrat, puis la
chambre de compensation unique qui fait de toutes ces bourses un seul marché, et la
soirée du troisième vendredi, où les options redeviennent des actions. Ce qu'une option
\emph{vaut} fait l'objet des livres~5 et~6~; ici la question est
de savoir ce qu'elle \emph{est}.

\section{Le contrat}

\begin{definition}[Call, put, prix d'exercice, échéance, prime]\label{def:m1:option-contracts-and-exchanges:option}
Une \emph{option d'achat}\index{option d'achat} confère à son détenteur le droit, et non
l'obligation d'acheter le sous-jacent à un \emph{prix d'exercice}\index{prix d'exercice} fixé, à
une date d'\emph{échéance}\index{echeance@échéance} au plus tard~; une \emph{option de vente}\index{option de vente} confère
le droit de vendre. Le prix payé pour l'option est sa
\emph{prime de l'option}\index{prime de l'option}. Le vendeur, ou émetteur, reçoit la
prime et supporte l'obligation.
\end{definition}

\begin{definition}[Multiplicateur, série, moneyness]\label{def:m1:option-contracts-and-exchanges:series}
Le \emph{multiplicateur d'option}\index{multiplicateur d'option} est la quantité de
sous-jacent par contrat~: 100 actions pour une option standard sur action américaine, 100~\$
par \omterm{def:m1:index-and-single-stock-futures-worldwide:point}{point d'indice} pour les grandes options sur indices. Une \emph{série d'options}\index{série d'options}
est l'ensemble des contrats partageant le même sous-jacent, le même type, le même
\omterm{def:m1:option-contracts-and-exchanges:option}{prix d'exercice} et la même \omterm{def:m1:option-contracts-and-exchanges:option}{échéance}~: l'unité qui porte un symbole, une cotation et
un carnet d'ordres. Un call est \emph{dans la monnaie}\index{dans la monnaie} lorsque le
sous-jacent est au-dessus de son \omterm{def:m1:option-contracts-and-exchanges:option}{prix d'exercice}, un put lorsqu'il est en dessous~; le montant est la
\emph{valeur intrinsèque}\index{valeur intrinsèque} de l'option.
\end{definition}

\begin{definition}[Exercice américain et exercice européen]\label{def:m1:option-contracts-and-exchanges:style}
Une option à \emph{exercice américain}\index{exercice américain} peut être
exercée à tout jour ouvré jusqu'à l'\omterm{def:m1:option-contracts-and-exchanges:option}{échéance}~; une option à
\emph{exercice européen}\index{exercice européen} ne peut l'être qu'à l'\omterm{def:m1:option-contracts-and-exchanges:option}{échéance}. Les options sur actions et
fonds américains sont à exercice américain et se dénouent par livraison des
actions~; les grandes options sur indices sont à exercice européen et se règlent en espèces.
\end{definition}

\begin{omfigure}
\begin{tikzpicture}
\begin{groupplot}[group style={group size=2 by 1,horizontal sep=1.6cm},
  width=6.4cm,height=4.6cm,
  xlabel={cours de l'action à l'échéance},
  tick label style={font=\small},label style={font=\small},
  xmin=40,xmax=60,ymin=-9,ymax=9,grid=major,grid style={gray!25},
  legend style={font=\footnotesize,at={(0.5,-0.36)},anchor=north,draw=none,fill=none,legend columns=2,/tikz/every even column/.append style={column sep=6pt}},
  table/col sep=comma]
\nextgroupplot[ylabel={P\&L par action (\$)},title={\small call, prix d'exercice 50, prime 2.40}]
\addplot[omDef,very thick] table[x=s,y=long_call]{figdata/markets-1/23-option-contracts-and-exchanges/payoffs.csv};
\addplot[omThm,very thick,dashed] table[x=s,y=short_call]{figdata/markets-1/23-option-contracts-and-exchanges/payoffs.csv};
\legend{acheteur,émetteur}
\nextgroupplot[title={\small put, prix d'exercice 50, prime 1.90}]
\addplot[omDef,very thick] table[x=s,y=long_put]{figdata/markets-1/23-option-contracts-and-exchanges/payoffs.csv};
\addplot[omThm,very thick,dashed] table[x=s,y=short_put]{figdata/markets-1/23-option-contracts-and-exchanges/payoffs.csv};
\legend{acheteur,émetteur}
\end{groupplot}
\end{tikzpicture}

\omcaption{P\&L à l'\omterm{def:m1:option-contracts-and-exchanges:option}{échéance}, prime incluse. La perte de l'acheteur est limitée à
la prime~; le gain de l'émetteur l'est aussi. Les deux courbes de chaque panneau
ont une somme nulle~: les options se répartissent, elles ne créent rien. Données~: la
construction du chapitre.}\label{fig:m1:option-contracts-and-exchanges:payoffs}
\end{omfigure}

Le symbole d'une option cotée aux États-Unis épelle entièrement sa série. Depuis 2010, chaque option
cotée porte une clé de 21 caractères~: six caractères pour la racine, complétés par
des espaces~; l'\omterm{def:m1:option-contracts-and-exchanges:option}{échéance} sur six chiffres, année, mois, jour~; \texttt{C} ou \texttt{P}~;
et le \omterm{def:m1:option-contracts-and-exchanges:option}{prix d'exercice} sur huit chiffres, cinq pour les dollars et trois pour les décimales.
\texttt{XYZ\textvisiblespace\textvisiblespace\textvisiblespace261218C00052500}
est le call XYZ qui expire le 18 décembre 2026, avec un \omterm{def:m1:option-contracts-and-exchanges:option}{prix d'exercice} de 52.50~\$.

Le nombre de séries cotées pour un titre découle des règles d'admission à la cote des bourses~:
des \omterm{def:m1:option-contracts-and-exchanges:option}{échéances} hebdomadaires pour les prochaines semaines, mensuelles pour les
prochains mois, trimestrielles et de long terme au-delà, chacune avec des \omterm{def:m1:option-contracts-and-exchanges:option}{prix
d'exercice} espacés selon un écart qui s'élargit à mesure qu'on s'éloigne de la monnaie et dans le temps. Cinq \omterm{def:m1:option-contracts-and-exchanges:option}{échéances} hebdomadaires
et six mensuelles, à soixante-dix \omterm{def:m1:option-contracts-and-exchanges:option}{prix d'exercice} chacune, et cinq \omterm{def:m1:option-contracts-and-exchanges:option}{échéances} lointaines avec
trente \omterm{def:m1:option-contracts-and-exchanges:option}{prix d'exercice}, donnent $2 \times (11 \times 70 + 5 \times 30) = 1\,840$ séries sur une
même action.

\section{Une seule chambre de compensation}

Aux États-Unis, toute option cotée, quelle que soit la bourse sur laquelle elle est négociée,
est émise et garantie par une seule chambre de compensation détenue par les bourses.
Le contrat de l'acheteur est conclu avec la chambre de compensation, et non avec l'émetteur, comme pour
toute \omterm{def:m1:clearing-and-settlement:ccp}{contrepartie centrale} (\cref{def:m1:clearing-and-settlement:ccp})~; ce qui est
inhabituel, c'est qu'il n'y en a qu'une. Un call acheté sur une bourse peut être revendu
sur n'importe quelle autre, car c'est le même contrat avec la même contrepartie~:
les options sont \emph{fongibles} d'une bourse à l'autre, ce qui n'est pas le cas des
contrats à terme (\cref{ch:m1:futures-contracts-and-exchanges}).

\begin{dated}{2026-09}{Taille}\label{dat:m1:option-contracts-and-exchanges:occ}
La chambre de compensation a publié 10.38 milliards de contrats compensés en 2022, puis
un nouveau record. Elle \omterm{def:m1:option-contracts-and-exchanges:assignment}{assigne} au hasard les avis d'exercice à ses membres compensateurs~; chaque
membre les répartit ensuite entre ses clients selon sa propre procédure.
\end{dated}

\section{Les bourses}

La fongibilité a rendu possible la concurrence entre bourses, et le résultat
ressemble à celui du marché des actions~: plus de quinze bourses d'options, détenues par une
poignée de groupes, toutes cotant les mêmes séries~; un meilleur cours acheteur et
un meilleur cours vendeur consolidés sur l'ensemble~; une règle interdisant de traiter en contournant un meilleur prix affiché~;
et une famille de modèles tarifaires, des bourses maker-taker à \omterm{def:m1:matching-algorithms-and-implied-spreads:fifo}{priorité prix-temps}
jusqu'aux bourses qui font payer le teneur, rémunèrent les ordres clients et
répartissent au prorata avec des privilèges pour les \omterm{def:m1:what-a-trading-firm-does:market-maker}{teneurs de marché} désignés
(\cref{ch:m1:matching-algorithms-and-implied-spreads}). Les options sur indices font
exception~: les plus importantes sont concédées en exclusivité à un seul groupe de
bourses et ne se négocient que là. La structure, et le flux qui la traverse, font
l'objet du \cref{ch:m1:options-market-structure}.

En Europe, la situation est celle des contrats à terme. Chaque bourse cote des options sur les
actions et les indices de son marché domestique et les compense dans sa propre chambre de
compensation~; une position ouverte sur une bourse ne peut pas être clôturée sur une autre, et la
liquidité de chaque sous-jacent se concentre en un seul endroit. Une part importante des volumes
européens en grande taille se négocie hors carnet d'ordres et se déclare à la bourse sous forme
de bloc.

\section{Exercice, assignation et échéance}

\begin{definition}[Assignation]\label{def:m1:option-contracts-and-exchanges:assignment}
Lorsqu'un porteur exerce, la chambre de compensation désigne un membre compensateur
ayant une position vendeuse sur la série et lui \emph{assigne}\index{assignation} l'obligation~; ce
membre transmet l'assignation à l'un de ses clients qui est vendeur. Le vendeur assigné d'un call
doit livrer les actions au \omterm{def:m1:option-contracts-and-exchanges:option}{prix d'exercice}~; le vendeur assigné d'un put doit les
acheter. Un vendeur apprend
l'assignation le lendemain matin.
\end{definition}

\begin{dated}{2026-09}{Les règles de l'échéance}\label{dat:m1:option-contracts-and-exchanges:expiry}
\emph{Exercice par exception.} À l'\omterm{def:m1:option-contracts-and-exchanges:option}{échéance}, la chambre de compensation exerce toute option
\omterm{def:m1:option-contracts-and-exchanges:series}{dans la monnaie} d'au moins 0.01~\$, qu'il s'agisse d'options sur actions ou sur indices, sauf
instruction contraire du membre compensateur. Il s'agit d'une commodité entre la chambre de
compensation et ses membres, et les courtiers peuvent appliquer leurs propres seuils.
\emph{Instructions.} Les porteurs peuvent aussi exercer des options qui ne sont pas \omterm{def:m1:option-contracts-and-exchanges:series}{dans la monnaie},
ou renoncer à exercer celles qui le sont~: la limite des bourses pour les notifications d'exercice
est fixée à 16:30, heure de Chicago, quatre-vingt-dix minutes après la clôture, et les heures limites
de la plupart des courtiers sont antérieures. \emph{Options sur indices.} Les options standard
sur l'indice S\&P 500 cessent de se négocier le jeudi et se règlent le troisième vendredi
sur la base d'une valeur calculée à partir des cours d'ouverture des composantes~; les séries
hebdomadaires et de fin de mois se négocient jusqu'à 16:00, heure de New York, à leur date
d'\omterm{def:m1:option-contracts-and-exchanges:option}{échéance} et se règlent sur les cours de clôture des composantes. Les deux sont européennes et
versent en espèces le montant \omterm{def:m1:option-contracts-and-exchanges:series}{dans la monnaie} multiplié par 100~\$.
\end{dated}

\begin{omfigure}
\begin{tikzpicture}[font=\small,
  box/.style={draw=omDef,rounded corners=2pt,minimum width=2.5cm,minimum height=1.0cm,align=center,fill=omDef!4},
  flow/.style={-{Stealth[length=2mm]},thick}]
\node[box] (h) at (0,0) {Porteur};
\node[box] (b1) at (3.3,0) {Courtier du\\porteur};
\node[box,draw=omThm,fill=omThm!5] (c) at (6.6,0) {Chambre de\\compensation};
\node[box] (b2) at (9.9,0) {Un membre\\vendeur};
\node[box] (w) at (13.0,0) {Un vendeur};
\draw[flow,omDef] (h) -- (b1);
\draw[flow,omDef] (b1) -- (c);
\draw[flow,omThm] (c) -- (b2);
\draw[flow,omThm] (b2) -- (w);
\node[font=\footnotesize,text=omDef] at (3.3,1.0) {notification d'exercice, avant la limite};
\node[font=\footnotesize,text=omThm] at (9.9,1.0) {assignation, pendant la nuit};
\node[font=\footnotesize] at (6.6,-1.0) {aléatoire};
\node[font=\footnotesize] at (9.9,-1.0) {règle propre au membre};
\end{tikzpicture}

\omcaption{Exercice et \omterm{def:m1:option-contracts-and-exchanges:assignment}{assignation}. Le porteur décide après la clôture~; le
vendeur l'apprend le lendemain matin. Entre les deux, le vendeur ne connaît
pas sa propre position.}\label{fig:m1:option-contracts-and-exchanges:assign}
\end{omfigure}

\begin{definition}[Risque de pin]\label{def:m1:option-contracts-and-exchanges:pin}
Le \emph{risque de pin}\index{risque de pin} est l'incertitude d'un vendeur dont les options vendues
expirent avec le sous-jacent à, ou très près de, le \omterm{def:m1:option-contracts-and-exchanges:option}{prix d'exercice}~: il ne peut pas savoir
combien seront exercées, donc combien d'actions il détiendra le lundi, et il ne peut pas couvrir
ce qu'il ignore.
\end{definition}

L'exercice est une décision prise par les porteurs avec quatre-vingt-dix minutes
de recul. Une option hors de la monnaie de 0.02~\$ à 16:00 est \omterm{def:m1:option-contracts-and-exchanges:series}{dans la monnaie} si
l'action se négocie 0.03~\$ plus haut dans la séance après clôture, à la suite d'une annonce
publiée à 16:05, et son porteur peut l'exercer. La \Cref{fig:m1:option-contracts-and-exchanges:expiry}
montre ce qu'une \omterm{def:m1:option-contracts-and-exchanges:option}{échéance} fait d'un petit carnet, en fonction du seul cours de clôture~;
les sauts se situent aux \omterm{def:m1:option-contracts-and-exchanges:option}{prix d'exercice}, et à chaque saut la réponse exacte est une
variable aléatoire.

\begin{omfigure}
\begin{tikzpicture}
\begin{axis}[width=11cm,height=5cm,
  xlabel={cours de clôture le jour de l'échéance},ylabel={actions nettes},
  tick label style={font=\small},label style={font=\small},
  xmin=46,xmax=54.5,ymin=-5200,ymax=7200,grid=major,grid style={gray!25},
  scaled y ticks=false,yticklabel style={/pgf/number format/fixed,/pgf/number format/1000 sep={\,}},
  table/col sep=comma]
\addplot[omDef,very thick,const plot] table[x=close,y=shares]{figdata/markets-1/23-option-contracts-and-exchanges/expiry.csv};
\node[font=\footnotesize,anchor=west] at (axis cs:46.05,5900) {puts vendus assignés};
\node[font=\footnotesize,anchor=west] at (axis cs:50.1,3400) {calls achetés exercés};
\node[font=\footnotesize,anchor=east] at (axis cs:54.45,-4350) {calls vendus assignés};
\end{axis}
\end{tikzpicture}

\omcaption{Un carnet vendeur de 40 puts à 47.50 et de 10 à 50, acheteur de 25 calls à 50
et vendeur de 60 calls à 52.50, tous arrivant à \omterm{def:m1:option-contracts-and-exchanges:option}{échéance} aujourd'hui. Actions reçues ($+$) ou livrées ($-$)
selon l'exercice par exception, chaque vendeur \omterm{def:m1:option-contracts-and-exchanges:series}{dans la monnaie} étant assigné en totalité,
en fonction du cours de clôture. Données~: la construction
du chapitre.}\label{fig:m1:option-contracts-and-exchanges:expiry}
\end{omfigure}

Les exercices anticipés d'options américaines sont rationnels dans deux cas qu'un desk
doit anticiper chaque jour, et pas seulement à l'\omterm{def:m1:option-contracts-and-exchanges:option}{échéance}~: un call juste avant une \omterm{def:m1:shares-corporate-actions-indices:dividend}{date
de détachement}, lorsque le dividende dépasse la valeur temps résiduelle de l'option~; et un put
profondément \omterm{def:m1:option-contracts-and-exchanges:series}{dans la monnaie}, lorsque l'intérêt sur le \omterm{def:m1:option-contracts-and-exchanges:option}{prix d'exercice} le dépasse. Le vendeur
de telles options est assigné pendant la nuit et se réveille avec une position en actions
et sans le dividende.

\section{Tutoriel~: symboles, une chaîne, profils de gains}\label{tut:m1:option-contracts-and-exchanges:chain}

\begin{tutorial}
\textbf{Objectif.} Analyser et produire des symboles OSI, conserver une chaîne dans une structure
indexée par des entiers, calculer les profils de gains et transformer un carnet en actions à l'\omterm{def:m1:option-contracts-and-exchanges:option}{échéance}.
\textbf{État final~:} les deux figures de données de ce chapitre.
\begin{tutsteps}
\item \textbf{Une série} avec son symbole, sa \omterm{def:m1:option-contracts-and-exchanges:series}{valeur intrinsèque} et son profil de gain.
\omcode{code/firm/chain/firm_chain.py}{11}{35}{Une série d'options. Le prix d'exercice est un entier de millièmes de dollar.}\label{lst:m1:option-contracts-and-exchanges:series}
\item \textbf{Analyse.} Vingt et un caractères, quatre champs, et une erreur pour
tout le reste.
\omcode{code/firm/chain/firm_chain.py}{38}{45}{D'un symbole OSI à une série.}\label{lst:m1:option-contracts-and-exchanges:parse}
\item \textbf{\omterm{def:m1:option-contracts-and-exchanges:option}{Échéance}.}
\omcode{code/firm/chain/firm_chain.py}{82}{98}{Actions reçues ou livrées à l'échéance selon l'exercice par exception.}\label{lst:m1:option-contracts-and-exchanges:expiry}
\end{tutsteps}
\textbf{Pour aller plus loin.} Rendre l'\omterm{def:m1:option-contracts-and-exchanges:assignment}{assignation} aléatoire~: chaque contrat vendeur situé à
moins de 0.05~\$ de la monnaie est exercé contre avec une probabilité qui dépend de la distance
et d'un mouvement en séance après clôture. Présenter la distribution de la position en actions
du lundi, et non un nombre.
\end{tutorial}

\section{Construction~: la chaîne d'options}\label{bld:m1:option-contracts-and-exchanges:chain}

\begin{build}
\bfield{Rôle} Tout composant de la mini-firme consacré aux options, de la surface de
volatilité du livre~5 au \omterm{def:m1:what-a-trading-firm-does:market-maker}{teneur de marché} du livre~11, demande
à un seul conteneur la série d'un sous-jacent. \bfield{Interface} \texttt{Series(root, expiry, right, strike\_milli,
multiplier, style, settlement)} avec \texttt{osi}, \texttt{strike},
\texttt{intrinsic}, \texttt{payoff}; \texttt{parse\_osi(symbol)};
\texttt{Chain(series)} avec \texttt{expiries()}, \texttt{strikes(expiry)}, \texttt{get(expiry, strike\_milli, right)}, \texttt{atm\_strike(expiry,
underlying)}~; \texttt{exercised\_by\_exception(series, close)};
\texttt{expiry\_shares(positions, close, overrides)}.
\bfield{Règles} Les strikes sont partout des entiers exprimés en millièmes~; un strike
flottant n'est jamais une clé. Les séries en double sont une erreur. Le strike à la
monnaie départage les égalités vers le bas, de façon déterministe. Les séries \omterm{def:m1:basis-roll-and-delivery:delivery}{réglées en espèces}
Les opérations sur titres qui modifient l'actif livrable
(\cref{bld:m1:shares-corporate-actions-indices:adjuster}) créent des séries ajustées
aux multiplicateurs non standard~: le multiplicateur est un champ, et non une
constante.
\bfield{Tests d'acceptation} \texttt{code/firm/chain/tests/}~: allers-retours de symboles
et symboles mal formés~; recherche dans la chaîne et règle du strike à la monnaie~; payoffs
des positions longues et courtes~; exercice par exception à un centime~; positions nettes
en actions d'un livre à plusieurs cours de clôture, avec une instruction contraire.
\bfield{Pour aller plus loin} Séries ajustées après un split et un dividende exceptionnel, avec
l'actif livrable sous forme de panier.
\end{build}

\begin{omsources}
\item The Options Industry Council, \emph{Options Exercise} (questions fréquentes sur
l'exercice~: seuils, heure limite, assignation).
\item The Options Clearing Corporation, communiqué de presse sur le volume compensé
en 2022, 3 janvier 2023~; documentation \emph{Options Symbology Initiative}.
\item Cboe, \emph{S\&P 500 Index Options: product specifications} et
\emph{S\&P 500 Weeklys Options: specifications}.
\end{omsources}

\section{Exercices}

\begin{exercise}[$\star$]\label{exo:m1:option-contracts-and-exchanges:1}
Décoder \texttt{XYZ\textvisiblespace\textvisiblespace\textvisiblespace270115P00047500}
et écrire le symbole du call XYZ du 19 mars 2027 avec un strike de
112.50~\$.
\end{exercise}

\begin{exercise}[$\star$]\label{exo:m1:option-contracts-and-exchanges:2}
Un trader achète 10 calls avec un strike de 50 pour une prime de 2.40. Donner le
P\&L à l'\omterm{def:m1:option-contracts-and-exchanges:option}{échéance} pour une action à 56 et à 48, ainsi que le cours de l'action au
seuil de rentabilité.
\end{exercise}

\begin{exercise}[$\star$]\label{exo:m1:option-contracts-and-exchanges:3}
Avec cinq \omterm{def:m1:option-contracts-and-exchanges:option}{échéances} hebdomadaires, six \omterm{def:m1:option-contracts-and-exchanges:option}{échéances} mensuelles de soixante-dix strikes chacune et cinq
\omterm{def:m1:option-contracts-and-exchanges:option}{échéances} lointaines de trente, dénombrer les séries. Combien de mises à jour de
cotations un mouvement d'un centime de l'action exige-t-il d'un \omterm{def:m1:what-a-trading-firm-does:market-maker}{teneur de marché} qui cote toutes
les séries sur seize bourses~?
\end{exercise}

\begin{exercise}[$\star\star$]\label{exo:m1:option-contracts-and-exchanges:4}
Pour le livre de la \cref{fig:m1:option-contracts-and-exchanges:expiry}, donner les
actions reçues ou livrées pour des cours de clôture de 47.00, 49.00, 50.00,
50.01, 52.50 et 52.51.
\end{exercise}

\begin{exercise}[$\star\star$]\label{exo:m1:option-contracts-and-exchanges:5}
Un call de strike 40 conserve encore 0.15 de valeur temps~; l'action se traite à
46 et passe ex-dividende demain, avec un dividende de 0.60. Le porteur doit-il
exercer aujourd'hui~? Que se passe-t-il pour le vendeur qui a couvert sa position avec des actions et qui est
assigné~?
\end{exercise}

\begin{exercise}[$\star\star$]\label{exo:m1:option-contracts-and-exchanges:6}
Une \omterm{def:m1:option-contracts-and-exchanges:series}{série d'options} sur indice cesse d'être négociée le jeudi et se règle le vendredi
sur les cours d'ouverture~; une autre se négocie jusqu'à la clôture du vendredi et se règle
sur les cours de clôture. Un desk est vendeur des deux, avec le même strike. Décrire
son risque entre la clôture du jeudi et la clôture du vendredi dans chacun des cas.
\end{exercise}

\begin{exercise}[$\star\star\star$]\label{exo:m1:option-contracts-and-exchanges:7}
\emph{Programmation.} Avec \texttt{expiry\_shares} et le livre du chapitre, ajouter une
instruction contraire~: le porteur des 50 calls longs refuse d'exercer à un cours de
clôture de 50.01. Donner les actions nettes avec et sans elle, et expliquer quand
refuser est justifié.
\end{exercise}

\begin{exercise}[$\star\star\star$]\label{exo:m1:option-contracts-and-exchanges:8}
\emph{Trouver la faille.} «~Le titre a clôturé à 52.48. Mes calls vendeurs à 52.50
ont expiré sans valeur, je garde la prime, et ma couverture peut être levée à l'ouverture du
lundi.~» Qu'a supposé le vendeur, et qu'aurait-il dû faire à
15:59~?
\end{exercise}

\section{Problème~: vendredi d'échéance}

\begin{problem}[{Problème du week-end --- la situation du livre lundi}]\label{pb:m1:option-contracts-and-exchanges:1}
Un desk détient le livre de la \cref{fig:m1:option-contracts-and-exchanges:expiry}
sur XYZ, toutes arrivant à \omterm{def:m1:option-contracts-and-exchanges:option}{échéance} aujourd'hui~: vendeur de 40 puts à 47.50, vendeur de 10 puts à 50, acheteur de
25 calls à 50, vendeur de 60 calls à 52.50. Multiplicateur 100. Il est couvert par une
position vendeuse de 1\,900 actions. XYZ clôture à 52.48.

\medskip\noindent\textbf{Partie I --- À la clôture.}
\begin{enumerate}
\item Quelles séries sont \omterm{def:m1:option-contracts-and-exchanges:series}{dans la monnaie} à 52.48, et de combien~?
\item Dans le cadre de l'exercice par exception, que devient chaque position~?
\item Donner les actions reçues à l'exercice.
\item Donner la position en actions du desk le lundi si rien d'autre ne se produit.
\item Que doit faire le desk de cette position, et quand~?
\end{enumerate}

\medskip\noindent\textbf{Partie II --- Après la clôture.}
À 16:10, XYZ annonce un contrat et se traite à 53.20 pendant la
séance après clôture.
\begin{enumerate}[resume]
\item Jusqu'à quand les porteurs des calls à 52.50 peuvent-ils transmettre leurs
instructions d'exercice~?
\item Pourquoi un porteur exercerait-il un call qui a clôturé hors de la monnaie~?
\item Le samedi, le desk apprend que 45 de ses 60 calls vendus ont été
assignés. Donner les actions livrées.
\item Donner la position en actions du desk le lundi.
\item XYZ ouvre le lundi à 54.10. Donner la perte sur la partie imprévue de
la position, par rapport au strike.
\end{enumerate}

\medskip\noindent\textbf{Partie III --- Ce qui aurait pu être fait.}
\begin{enumerate}[resume]
\item À 15~h~55, XYZ à 52.45, les calls à 52.50 étaient offerts à 0.04. Combien
aurait coûté le rachat des 60~?
\item Comparer avec la question 10.
\item Pourquoi l'option coûte-t-elle encore 0.04 cinq minutes avant
d'«~expirer sans valeur~»~?
\item Le desk aurait pu acheter à la place 3\,000 actions à la clôture, par
précaution. Évaluer.
\item Pourquoi le long call à 50 n'est-il pas une préoccupation de même nature~?
\end{enumerate}

\medskip\noindent\textbf{Partie IV --- Jugement.}
\begin{enumerate}[resume]
\item On dit que la couverture des \omterm{def:m1:what-a-trading-firm-does:market-maker}{teneurs de marché} «~épingle~» les actions aux
strikes le jour de l'\omterm{def:m1:option-contracts-and-exchanges:option}{échéance}. Donner le mécanisme en deux phrases.
\item Ce problème existerait-il pour une option européenne sur indice réglée en espèces~?
\item Un client de détail est assigné sur un call vendu, dans une action qui
ouvre ensuite en forte hausse. Où se situe la perte jusqu'au lundi~?
\item Énoncer le \emph{résultat nommé}~: la position en actions du desk le lundi
après exercice et \omterm{def:m1:option-contracts-and-exchanges:assignment}{assignation}.
\item En une phrase~: qui décide du moment où une position en options devient une
position en actions~?
\end{enumerate}
\end{problem}

\section{Questions d'entretien}

\begin{interviewq}[$\star$ \iqroles{trader, chercheur, développeur}]\label{iq:m1:option-contracts-and-exchanges:1}
Quelle est la différence entre les options américaines et européennes, et
quels produits cotés relèvent de chacune~?
\end{interviewq}

\begin{interviewq}[$\star$ \iqroles{trader, développeur}]\label{iq:m1:option-contracts-and-exchanges:2}
Pourquoi une option sur action des États-Unis achetée sur une bourse peut-elle être
vendue sur une autre, alors qu'un contrat à terme ne le peut pas~?
\end{interviewq}

\begin{interviewq}[$\star\star$ \iqroles{trader, chercheur}]\label{iq:m1:option-contracts-and-exchanges:3}
Quand est-il rationnel d'exercer un call américain par anticipation~? Et un put~?
\end{interviewq}

\begin{interviewq}[$\star\star$ \iqroles{trader}]\label{iq:m1:option-contracts-and-exchanges:4}
Qu'est-ce que le \omterm{def:m1:option-contracts-and-exchanges:pin}{risque de pin} et comment le gérez-vous~?
\end{interviewq}

\begin{interviewq}[$\star\star$ \iqroles{développeur, chercheur}]\label{iq:m1:option-contracts-and-exchanges:5}
Concevez la clé d'un instrument optionnel dans un système de trading.
\end{interviewq}

\begin{interviewq}[$\star\star\star$ \iqroles{trader, chercheur}]\label{iq:m1:option-contracts-and-exchanges:6}
Vous êtes vendeur d'une option sur indice qui se règle sur les cours d'ouverture
du matin de l'\omterm{def:m1:option-contracts-and-exchanges:option}{échéance}. Comment couvrez-vous la nuit précédente~?
\end{interviewq}
